\chapter{Weihrauch Reducibility, Choice and Omniscience}
\label{sec:p04-weihrauch-choice}

A computational problem may permit several correct answers. Its uniform
strength concerns the information needed to produce an allowed answer on
every promised input name, not whether some answer exists classically.
This chapter extends the represented-space foundation of
\cref{sec:p04-represented-spaces} to such problems. The presentations used
are Brattka--Gherardi's July 2009 author version
\cite{BrattkaGherardi2009OmniscienceV2} and
Brattka--de Brecht--Pauly's September 2010 author version
\cite{BrattkaDeBrechtPauly2010ChoiceV2}.

\section{Realizers and the two reduction interfaces}
\label{sec:p04-wc-realizers}

\begin{hbdefinition}{def:p04-wc-multivalued}{A multivalued problem and all-name realization}
A partial multivalued problem \(f:\subseteq X\rightrightarrows Y\) assigns a
nonempty set \(f(x)\subseteq Y\) to every \(x\in\dom(f)\).
Its domain is a promise; inputs outside it have no required answer.
For represented spaces \((X,\delta_X),(Y,\delta_Y)\), a single-valued
partial stream operator \(F:\mathcal B\rightharpoonup\mathcal B\)
is a realizer, written \(F\vdash f\), if on every name \(p\) of
every promised input it is defined, gives a valid output name, and
\[
 \delta_Y(F(p))\in f(\delta_X(p)).
\]
Different input names may lead to different allowed answers.
Neither a realizer nor a reduction may rely on a favourable selection
by a multivalued oracle. The quantifier over oracle realizers below
includes noncomputable realizers and every legitimate output-name choice.
\end{hbdefinition}

\begin{hbconvention}{def:p04-wc-tupling}{Name products and stream coordinates}
The binary pairing of streams interleaves their entries:
\(\langle p,q\rangle(2n)=p(n)\) and
\(\langle p,q\rangle(2n+1)=q(n)\).
For countable tuples, use the fixed effective numerical pairing to put
\(\langle p_0,p_1,\ldots\rangle(\langle i,j\rangle)=p_i(j)\).
All projections are computable. A represented sequence has these names
with its component representations; a natural answer is represented by
\(\delta_{\mathbb N}(q)=q(0)\), with arbitrary tail entries.
Binary streams have the identity representation on
\(\mathcal C:=\{0,1\}^{\mathbb N}\), the Cantor space.
\end{hbconvention}

\begin{hbdefinition}{def:p04-wc-reduction}{Ordinary and strong Weihrauch reduction}
In Brattka--Gherardi's notation, \(f\leq_{\mathrm W} g\) means that there
are computable partial stream operators \(H,K\) such that
\[
 H\langle\mathrm{id},GK\rangle\vdash f
                         \qquad\text{for every }G\vdash g.
\]
Here \(K\) preprocesses the input name and \(H\) receives both
that original name and the oracle's output. In particular \(K(p)\)
is a valid promised \(g\)-input name whenever \(p\) names a
promised \(f\)-input. Strong reduction, \(f\leq_{\mathrm{sW}}g\), instead
requires \(HGK\vdash f\) for every \(G\vdash g\): the
postprocessor does not receive the original input.
Write \(\equiv_{\mathrm W},\equiv_{\mathrm{sW}}\) for mutual reduction,
\(<_{\mathrm W}\) for reduction without the reverse reduction, and
\(\mid_{\mathrm W}\) for incomparability.

The 2010 closed-choice paper uses the distinct variable assignment
\(K\langle\mathrm{id},GH\rangle\vdash f\): its \(H\) is the
preprocessor and its \(K\) the postprocessor. This renaming of
the two operators does not change the interface; their roles are not
silently read in the opposite order.
\end{hbdefinition}

\begin{hbtheorem}{thm:p04-wc-basic}{proved=annex:p04-wc-basic}{Preorders, preservation and representation invariance}
Both reductions are reflexive and transitive. Strong reduction implies
ordinary reduction. If \(f\leq_{\mathrm W} g\) and \(g\) has a computable
or continuous realizer, then \(f\) has a realizer of that same kind.
Computably equivalent replacements of all input/output representations
preserve either reduction in both directions.
For the identity on Baire space and a total constant problem
\(c:\mathcal B\to\{0\}\), one has
\(\mathrm{id}\leq_{\mathrm W} c\) but
\(\mathrm{id}\not\leq_{\mathrm{sW}}c\).
\end{hbtheorem}

The ordinary equivalence classes are Weihrauch degrees. The degree
comparison is about one oracle call after computable preprocessing.
It does not automatically permit adaptive repeated oracle calls.
The constant example shows why losing the original input is a real
restriction even when neither problem requires a noncomputable answer.

\begin{hbdefinition}{def:p04-wc-cylinder}{Products and cylindrification}
The product problem has
\((f\times g)(x,u)=f(x)\times g(u)\), with domain
\(\dom(f)\times\dom(g)\), using product names.
The cylindrification of \(g\) is \(\mathrm{id}_{\mathcal B}\times g\).
A problem is a cylinder if
\(g\equiv_{\mathrm{sW}}\mathrm{id}_{\mathcal B}\times g\).
\end{hbdefinition}

\begin{hbtheorem}{thm:p04-wc-cylinder}{proved=annex:p04-wc-cylinder}{Exactly where ordinary and strong reduction agree}
For all problems,
\[
 f\leq_{\mathrm W}g\ \Longleftrightarrow\
 \mathrm{id}\times f\leq_{\mathrm{sW}}\mathrm{id}\times g.
\]
If \(g\) is a cylinder, then
\(f\leq_{\mathrm W} g\) if and only if \(f\leq_{\mathrm{sW}}g\).
\end{hbtheorem}

\section{Finite and countable parallelization}
\label{sec:p04-wc-parallel}

\begin{hbdefinition}{def:p04-wc-parallel}{Independent lists of instances}
The countable parallelization is
\[
 \widehat f((x_i)_{i\in\mathbb N})=
                 \prod_{i\in\mathbb N}f(x_i),
 \qquad
 \dom(\widehat f)=\dom(f)^{\mathbb N}.
\]
Each coordinate is an independently supplied instance; the output
contains an allowed answer at every coordinate. Countable tuple names
are as in \cref{def:p04-wc-tupling}.
Finite parallelization \(f^*\) takes a length-tagged finite list
\((n,x_0,\ldots,x_{n-1})\) and returns a list of the same length
with one allowed answer per instance. The zero-length list is a valid
input with the empty list as its sole decoded answer. Its zero-arity
constant problem is ordinarily, not strongly, equivalent to the identity.
\end{hbdefinition}

\begin{hbtheorem}{thm:p04-wc-parallel}{proved=annex:p04-wc-parallel}{Countable closure and tuple rearrangements}
For either reduction,
\[
 f\leq_{\mathrm{sW}}\widehat f,\qquad
 f\leq_{\mathrm W}g\ \Longrightarrow\ \widehat f\leq_{\mathrm W}\widehat g,
 \qquad \widehat{\widehat f}\equiv_{\mathrm{sW}}\widehat f.
\]
The monotonicity implication also holds with strong reduction throughout.
Furthermore,
\[
 \widehat{f\times g}\equiv_{\mathrm{sW}}\widehat f\times\widehat g,
 \qquad \widehat f\times\widehat f\equiv_{\mathrm{sW}}\widehat f.
\]
No computable point in \(\dom(f)\) is needed: duplication or
interleaving uses the supplied instances. The statements also cover an
empty promise domain.
\end{hbtheorem}

\begin{hbtheorem}{thm:p04-wc-finite}{proved=annex:p04-wc-finite}{Finite closure with the partition retained}
Finite parallelization is an ordinary Weihrauch closure operation:
\(f\leq_{\mathrm W}f^*\), monotonicity holds, and
\((f^*)^*\equiv_{\mathrm W}f^*\).
For a problem with a computable promised input,
\(f^*\leq_{\mathrm W}\widehat f\).
In the idempotence construction the ordinary postprocessor can recover
the original partition into sublists from the input. No general strong
idempotence assertion is inferred by discarding that partition.
\end{hbtheorem}

\section{Negative information and the omniscience interfaces}
\label{sec:p04-wc-choice-definitions}

\begin{hbdefinition}{def:p04-wc-negative}{Closed sets described by exclusions}
For a discrete \(X\subseteq\mathbb N\), a negative name of
\(A\subseteq X\) lists the excluded points plus one, with zero
padding, arbitrary order and repetitions:
\[
 \psi_-^X(p)=\{n\in X:n+1\notin\operatorname{range}(p)\}.
\]
Positive entries must denote points of \(X\). For \(X=\mathbb N\)
this is the 2010 paper's \(\psi_-^{\mathbb N}\) clause.
For Cantor space, fix the length-lexicographic enumeration
\((w_j)\) of finite binary words. A positive name entry \(j+1\)
excludes the cylinder
\([w_j]=\{b\in\mathcal C:w_j\text{ is a prefix of }b\}\);
zero is padding. Thus
\[
 \psi_-^{\mathcal C}(p)=
 \mathcal C\setminus\bigcup_{p(t)>0}[w_{p(t)-1}].
\]
This explicit coding implements the source's potentially empty
enumeration of exterior cylinders. It supplies negative information,
not a membership decision for the remaining closed set.
\end{hbdefinition}

\begin{hbdefinition}{def:p04-wc-choice}{Closed choice}
Closed choice is
\(\mathsf{C}_X:\subseteq\mathcal A_-(X)\rightrightarrows X\), with
\(\mathsf{C}_X(A)=A\) on the domain of nonempty negatively named closed sets.
Here \(\mathcal A_-(X)\) denotes the represented set of closed sets
with the specified negative representation. An output is a point of
\(A\), not a new name of the set \(A\).
We use \(\mathsf{C}_{\mathbb N}\), \(\mathsf{C}_{\{0,1\}}\), and
\(\mathsf{C}_{\mathcal C}=\mathsf{C}_{\{0,1\}^{\mathbb N}}\) below.
The empty-set input is outside the promise.
\end{hbdefinition}

\begin{hbdefinition}{def:p04-wc-omniscience}{The 2009 LPO and LLPO presentations}
Retaining Definition 6.1 of the 2009 author text,
\[
 \mathsf{LPO}(p)=
 \begin{cases}0,&(\exists n)\ p(n)=0,\\1,&(\forall n)\ p(n)\ne0,
 \end{cases}
 \qquad(p\in\mathcal B).
\]
Also
\[
 \mathsf{LLPO}(p)=\{i\in\{0,1\}:(\forall n)\ p(2n+i)=0\},
\]
on inputs with at most one nonzero entry. In particular,
\(\mathsf{LLPO}(0^{\mathbb N})=\{0,1\}\).
The alternative zero-stream test
\(\mathsf{LPO}_0(p)=1\) exactly when \(p=0^{\mathbb N}\)
is not the same literal map as the displayed \(\mathsf{LPO}\).
They are strongly equivalent: complement the coordinatewise zero/nonzero
indicator in the input. The annex verifies both translations.
\end{hbdefinition}

\begin{hbdefinition}{def:p04-wc-mlpo}{The 2010 tuple presentation of finite omniscience}
In the distinct 2010 presentation, for \(n\geq1\),
\[
 \mathsf{MLPO}_n:\subseteq\mathcal B\rightrightarrows\mathbb N,
 \qquad \mathsf{MLPO}_n\langle p_1,\ldots,p_n\rangle
       =\{i\in\{1,\ldots,n\}:p_i=0^{\mathbb N}\},
\]
with domain
\(\{\langle p_1,\ldots,p_n\rangle:(\exists i\in\{1,\ldots,n\})
\ p_i=0^{\mathbb N}\}\). Thus at least one component is the zero stream.
The source writes \(\mathsf{LLPO}=\mathsf{MLPO}_2\) for this
two-component problem. Its answer labels are one and two, and a
nonzero component may contain many nonzero entries. This is not the
same literal input convention or answer labeling as the 2009
single-marker \(\mathsf{LLPO}\) in \cref{def:p04-wc-omniscience}.
\end{hbdefinition}

\begin{hbtheorem}{thm:p04-wc-finite-choice}{proved=annex:p04-wc-finite-choice}{Finite closed choice and the tuple presentation}
For every fixed \(n\geq0\),
\[
 \mathsf{C}_{\{0,\ldots,n\}}\equiv_{\mathrm{sW}}
                                      \mathsf{MLPO}_{n+1}.
\]
In particular the 2010 tuple presentation and 2009 single-marker
presentation of \(\mathsf{LLPO}\) are strongly equivalent, by
explicit translations through binary closed choice. The shift from
zero-based point labels to one-based tuple answers is part of the
postprocessor.
\end{hbtheorem}

The parity-stream formulas and subsequent degree comparisons use the
2009 single-marker presentation of \(\mathsf{LLPO}\). The 2010 tuple
presentation keeps its distinct domain and labels in the definition
and finite-choice bridge above.

\begin{hbtheorem}{thm:p04-wc-binary-choice}{proved=annex:p04-wc-binary-choice}{Binary negative choice is LLPO}
\[
 \mathsf{C}_{\{0,1\}}\equiv_{\mathrm{sW}}\mathsf{LLPO}
 \leq_{\mathrm{sW}}\mathsf{LPO}\equiv_{\mathrm{sW}}\mathsf{LPO}_0.
\]
Both directions preserve the single-nonzero LLPO promise. An exclusion
can be repeated indefinitely in a negative name; the translator emits
one nonzero LLPO marker, not one marker for each repetition.
\end{hbtheorem}

\begin{hbtheorem}{thm:p04-wc-omniscience-obstruction}{proved=annex:p04-wc-omniscience-obstruction}{No continuous oracle-free realizer}
Neither \(\mathsf{LPO}\) nor \(\mathsf{LLPO}\) has a continuous
realizer, and hence neither has a computable one. The negative result
for LLPO holds despite allowing name-dependent choice at its two-answer
input. Consequently binary closed choice is not computable.
\end{hbtheorem}

\section{Three distinct limits and natural-number choice}
\label{sec:p04-wc-limits}

\begin{hbdefinition}{def:p04-wc-limit-types}{Coordinatewise, scalar discrete and whole-stream discrete limits}
The Baire limit \(\lim\) takes a sequence of streams \((p_s)\)
such that for every coordinate \(n\), \(p_s(n)\) is eventually
constant, and returns the stream of those eventual values. No stabilization
rate is supplied.
The scalar problem \(\lim_{\mathbb N}\) takes an eventually
constant sequence of natural numbers and returns its final value.
In contrast, the 2010 source's \(\lim_\Delta\) takes a sequence
of entire Baire streams that is eventually constant as a sequence of
whole streams, and returns that stream. Its input promise is stronger
than coordinatewise convergence. Its output can still be a noncomputable
stream. This is the limit for the discrete topology on Baire space,
not a limit with the discrete codomain \(\mathbb N\).
\end{hbdefinition}

\begin{hbtheorem}{thm:p04-wc-natural-limits}{proved=annex:p04-wc-natural-limits}{Discrete choices with and without input recovery}
\[
 \mathsf{C}_{\mathbb N}\equiv_{\mathrm{sW}}\lim_{\mathbb N},\qquad
 \mathsf{C}_{\mathbb N}\equiv_{\mathrm W}\lim_\Delta,
 \qquad \lim_\Delta\not\leq_{\mathrm{sW}} \mathsf{C}_{\mathbb N}.
\]
The ordinary reduction of \(\lim_\Delta\) chooses a stabilization
stage and then retrieves the original stream at that stage.
For the scalar limit one can instead encode both a stage and its
candidate value in the choice set, permitting strong postprocessing.
\end{hbtheorem}

\begin{hbtheorem}{thm:p04-wc-parallel-limit}{proved=annex:p04-wc-parallel-limit}{Countably many decisions and coordinatewise convergence}
\[
 \widehat{\mathsf{LPO}}\equiv_{\mathrm{sW}}\lim
              \equiv_{\mathrm{sW}}\widehat{\mathsf{C}_{\mathbb N}}.
\]
The reductions supply both directions with the actual name interfaces;
they do not treat a convergence promise as a known effective rate.
The ordinary discrete-limit equivalence in the preceding theorem does
not identify \(\lim_\Delta\) with these countable parallel problems.
\end{hbtheorem}

\section{Cantor choice, binary trees and parallel LLPO}
\label{sec:p04-wc-cantor}

\begin{hbdefinition}{def:p04-wc-wkl}{The characteristic-tree WKL interface}
A binary tree \(T\subseteq\{0,1\}^{<\mathbb N}\) is prefix-closed.
The source representation \(\delta_{\mathrm{Tr}}\) gives its
characteristic function on the fixed enumeration of binary words:
\(\delta_{\mathrm{Tr}}(p)=T\) when \(p(j)=\chi_T(w_j)\).
Let \([T]\) be the binary streams whose every finite prefix belongs
to \(T\). The problem \(\mathsf{WKL}(T)=[T]\) has the promise
that \(T\) is infinite. Classical binary-tree compactness guarantees
that this output set is nonempty; it does not compute a path.
\end{hbdefinition}

\begin{hbtheorem}{thm:p04-wc-cantor-tree}{proved=annex:p04-wc-cantor-tree}{Negative closed sets and characteristic trees}
\[
 \mathsf{C}_{\mathcal C}\equiv_{\mathrm{sW}}\mathsf{WKL}.
\]
From an exterior-cylinder enumeration one computes a decidable tree by
testing only the first \(n\) exclusions at depth \(n\).
The resulting path set is exactly the named closed set.
The reverse translator enumerates cylinders of nodes absent from the
characteristic tree.
\end{hbtheorem}

\begin{hbtheorem}{thm:p04-wc-cantor-llpo}{proved=annex:p04-wc-cantor-llpo}{A full pair of parallel choice constructions}
\[
 \mathsf{C}_{\mathcal C}\equiv_{\mathrm{sW}}\widehat{\mathsf{LLPO}}.
\]
One direction encodes coordinatewise admissible bits as a closed product
set. The other asks one LLPO question for every finite binary prefix
about which child cylinder is not empty. The queries are supplied
nonadaptively, and the answer table is used adaptively only by the
computable postprocessor constructing the path.
For an empty prefix-cylinder, both children may be empty: the query
generator must still emit at most one marker to remain a valid LLPO
instance. The proof supplies that guard.
\end{hbtheorem}

\begin{hbtheorem}{thm:p04-wc-omniscience-cylinders}{proved=annex:p04-wc-omniscience-cylinders}{Input recovery in the parallelized principles}
The problems \(\widehat{\mathsf{LPO}}\) and
\(\widehat{\mathsf{LLPO}}\) are cylinders. Therefore so are
\(\lim\), \(\widehat{\mathsf{C}_{\mathbb N}}\),
\(\mathsf{C}_{\mathcal C}\) and \(\mathsf{WKL}\), with the interfaces
specified here. Ordinary and strong reduction to any of these targets
coincide. The LLPO input-recovery code uses a single marker in each
query, respecting the source's domain restriction.
\end{hbtheorem}

\begin{hbdefinition}{def:p04-wc-real-llpo}{The source real comparison relation}
With the scalar Cauchy representation, define
\[
 \mathsf{LLPO}_{\mathbb R}(x)
 =\{0:x\leq0\}\cup\{1:x\geq0\}.
\]
At zero both answers are permitted. This is a total multivalued
weak comparison, not a total discrete sign decision with a separate
equality answer.
\end{hbdefinition}

\begin{hbtheorem}{thm:p04-wc-real-llpo}{proved=annex:p04-wc-real-llpo}{Real weak comparison and the single-marker problem}
\(\mathsf{LLPO}_{\mathbb R}\equiv_{\mathrm{sW}}\mathsf{LLPO}\)
for the 2009 single-marker presentation.
The corrected sign-marker translator puts a marker in the even
positions for a positive real and in the odd positions for a negative
real, and produces the zero stream when neither strict sign is witnessed.
\end{hbtheorem}

\section{Weak computability and exact separation boundaries}
\label{sec:p04-wc-boundaries}

\begin{hbdefinition}{def:p04-wc-weak}{Weakly computable problems}
A problem is weakly computable if
\(f\leq_{\mathrm W}\widehat{\mathsf{LLPO}}\), equivalently
\(f\leq_{\mathrm W} \mathsf{C}_{\mathcal C}\). This is a computational degree
definition, not a probabilistic confidence claim and not a synonym
for computable. The equivalence permits noncomputable multivalued
problems, including Cantor choice itself.
\end{hbdefinition}

\begin{hbtheorem}{thm:p04-wc-weak-single}{proved=annex:p04-wc-weak-single}{Single-valued Baire outputs cannot hide noncomputability in compact choice}
For a partial single-valued \(f:\mathcal B\rightharpoonup\mathcal B\),
using identity representations, if \(f\leq_{\mathrm W} \mathsf{C}_{\mathcal C}\)
then \(f\) is computable. The annex gives the direct finite-cover
construction. This is the Baire-output case of the source's more general
single-valued weak-computability result; arbitrary output representations
must not be assumed to give identical output-name digits.
\end{hbtheorem}

\begin{hbtheorem}{thm:p04-wc-no-computable-point}{proved=annex:p04-wc-no-computable-point}{An effective closed Cantor set without any computable element}
There is a nonempty co-computably-enumerable closed
\(A\subseteq\mathcal C\) with no computable binary stream.
Its negative name is computable. Consequently
\(\mathsf{C}_{\mathcal C}\not\leq_{\mathrm W} \mathsf{C}_{\mathbb N}\) and
\(\mathsf{C}_{\mathcal C}\not\leq_{\mathrm W}\mathsf{LPO}\).
An existence guarantee for a closed set does not guarantee a computable
point even when its negative description is computable.
\end{hbtheorem}

\begin{hbtheorem}{thm:p04-wc-degree-boundaries}{proved=annex:p04-wc-degree-boundaries}{The comparisons established by these constructions}
\[
 \mathsf{LLPO}<_{\mathrm W}\mathsf{LPO},\qquad
 \mathsf{LPO}\mid_{\mathrm W} \mathsf{C}_{\mathcal C},\qquad
 \mathsf{C}_{\mathbb N}\mid_{\mathrm W} \mathsf{C}_{\mathcal C},\qquad
 \mathsf{C}_{\mathcal C}<_{\mathrm W}\lim.
\]
Moreover \(\mathsf{LPO}\leq_{\mathrm{sW}}\mathsf{C}_{\mathbb N}\) and both
\(\mathsf{C}_{\mathbb N}\) and \(\mathsf{C}_{\mathcal C}\) reduce strongly to
\(\lim\). These statements do not identify finite LLPO products
with one LLPO call, finite with countable parallelization, or arbitrary
closed choice on Baire space with Cantor choice.
\end{hbtheorem}

For example, a negative name successively excluding \(0,1,2\) can
force a natural-choice approximation to change several times before it
stabilizes. A nonempty Cantor closed set instead permits a whole binary
path, and its computable descriptions can have no computable answer.
These are different failures of uniform computability. The exact
reductions and separation proofs, rather than their informal descriptions,
determine the comparisons above.
